\section{Ionization and chemical state}
\label{sec:ionization}

\subsection{The species carried}
\label{sec:ionization:species}

The composition of a cell is the vector \texttt{f\_sp}$(j,\cdot)$ of species
fractions, laid out once in \texttt{species\_table.f90}:

\begin{center}
\begin{tabular}{llp{0.46\textwidth}}
\toprule
column(s) & species & note\\
\midrule
1--2   & H\,{\sc i}, H\,{\sc ii} & fractions of the hydrogen nuclei\\
3--5   & He\,{\sc i}, He\,{\sc ii}, He\,{\sc iii} & fractions of the helium
         nuclei; column 3 \emph{contains} the metastable\\
6      & He\,$2^3$S & the metastable triplet, tracked when
         \texttt{Include He23S?} is true (default on)\\
7--33  & the 27 metal ion stages & C, O, N, Mg, Si, Ca, Fe in three stages;
         Na, K, S in two\\
34--37 & H$_2$, H$_2^+$, H$_3^+$, HeH$^+$ & molecular option
         (\texttt{thereis\_mol}, default off)\\
38--40 & OH, H$_2$O, CO & oxygen option (\texttt{thereis\_oxychem},
         default off)\\
\bottomrule
\end{tabular}
\end{center}

Ten elements carry a nuclear abundance $n_X/n_{\rm H}$, set by
\texttt{metals.inp} (\texttt{metals\_input\_read.f90}) or by a
lower-atmosphere handoff. Masses are in units of the hydrogen \emph{atom}:
helium weighs $m_{\rm He}/m_{\rm H} = 3.9715259$ and the metal weights are
the CIAAW values converted from the unified mass unit by
$1\,{\rm u} = 0.99223573\,m_{\rm H}$.

\subsection{The cell system}
\label{sec:ionization:cellsystem}

For each cell the code solves a system of algebraic balance equations at
fixed temperature, density and radiation field. The unknowns are stage
\emph{fractions}, and one stage of each element is closed as the remainder
rather than being an unknown: $n({\rm H\,I}) = (1-x_{\rm H^+})n_{\rm H}$,
$n({\rm He\,I}) = (1-x_{\rm He^+}-x_{\rm He^{2+}})n_{\rm He}$, and for a metal
element $n(X^{0}) = (1-x_{X^+}-x_{X^{2+}})n_X$. The layout is

\begin{center}
\begin{tabular}{ll}
\toprule
rows & contents\\
\midrule
1--3 & $x_{\rm H^+}$, $x_{\rm He^+}$, $x_{\rm He^{2+}}$\\
4    & He\,$2^3$S (triplet systems) \emph{or} H$_2$ (molecular systems)\\
5--8 & H$_2^+$, H$_3^+$, HeH$^+$, He\,$2^3$S (molecular systems)\\
$m_{\rm base}$\dots & two rows per metal element
       (\texttt{metal\_row\_base()})\\
$i_{\rm ox}$\dots & OH and H$_2$O (oxygen option)\\
\bottomrule
\end{tabular}
\end{center}

with $m_{\rm base} = 4$, or 5 with the triplet. In the molecular layout
$i_{\rm ox} = 8$, or 9 with the triplet, and
$m_{\rm base} = i_{\rm ox}$, raised by two when the oxygen chemistry is on
(\texttt{metal\_row\_base}, \texttt{oxygen\_row\_base}: one definition, read
by the residual and by the driver alike). Seven residual routines cover the
configurations: \texttt{System\_H}, \texttt{System\_HeH},
\texttt{System\_HeH\_metals}, \texttt{System\_HeH\_TR},
\texttt{System\_HeH\_TR\_metals}, \texttt{System\_HeH\_mol} and
\texttt{System\_HeH\_mol\_metals}. The shared rows live once, in
\texttt{ion\_residual\_core.f90}.

\paragraph{The H/He rows.} With photoionization rates $\Gamma$, case-B
recombination coefficients $\alpha$ and electron-impact ionization
coefficients $\beta$, \texttt{heh\_rows} writes
\begin{align}
0 &= n_{\rm HI}\Gamma_{\rm HI}
   + \left(n_{\rm HI}\beta_{\rm HI}-\alpha_{\rm HII}n_{\rm HII}\right)n_e,
   \label{eq:ion:hrow}\\
0 &= n_{\rm HeI}\Gamma_{\rm HeI}
   + \left(n_{\rm HeI}\beta_{\rm HeI}-\alpha_{\rm HeII}n_{\rm HeII}\right)n_e,
   \label{eq:ion:herow}\\
0 &= n_{\rm HeII}\Gamma_{\rm HeII}
   + \left(n_{\rm HeII}\beta_{\rm HeII}-\alpha_{\rm HeIII}n_{\rm HeIII}\right)n_e,
   \label{eq:ion:he2row}
\end{align}
with $n_e$ the charge sum of every tracked ion,
$n_e = n_{\rm HII}+n_{\rm HeII}+2n_{\rm HeIII}+\sum_i Z_i n_{{\rm m},i}
 + n_{\rm H_2^+}+n_{\rm H_3^+}+n_{\rm HeH^+}$
(\texttt{calc\_ne}), so charge neutrality is imposed identically rather
than solved for. Each metal element contributes
\begin{equation}
0 = n_{X^0}\Gamma_{X^0}
  + \left(n_{X^0}\beta_{X^0}-\alpha_{X^+}n_{X^+}\right)n_e,
\qquad
0 = n_{X^+}\Gamma_{X^+}
  + \left(n_{X^+}\beta_{X^+}-\alpha_{X^{2+}}n_{X^{2+}}\right)n_e,
\label{eq:ion:metalrows}
\end{equation}
the second row being replaced by $x_{X^{2+}}=0$ for a two-stage element, and
the whole pair by $x_{X^+}=x_{X^{2+}}=0$ for an element with zero abundance.

\paragraph{The triplet row.} With the metastable tracked, the He rows are
rewritten in the form of \citet{Oklopcic2018} (rows 2 and 3 summed, the triplet and the
singlet ground state distinguished) and the fourth row is the $2^3$S
balance,
\begin{equation}
0 = -n_{2^3{\rm S}}\Gamma_{2^3{\rm S}}
    + n_e\left[n_{\rm HeII}\alpha_{2^3{\rm S}}
             + n_{1^1{\rm S}}q_{13}
             - n_{2^3{\rm S}}(q_{31a}+q_{31b})
             - n_{2^3{\rm S}}\beta_{2^3{\rm S}}\right]
    - n_{2^3{\rm S}}\left(A_{31}+n_{\rm HI}Q_{31}\right),
\label{eq:ion:tripletrow}
\end{equation}
where $q_{13}$ is the collisional excitation out of the singlet ground
state, $q_{31a}$ and $q_{31b}$ the collisional conversions to $2^1$S and
$2^1$P, $A_{31} = 1.272\times10^{-4}$\,s$^{-1}$ the forbidden radiative
decay, $\beta_{2^3{\rm S}}$ the electron-impact ionization of the
metastable at its own 4.8\,eV threshold, and $Q_{31}$ the \emph{total}
He($2^3$S)$+$H ionization rate coefficient. Both channels of $Q_{31}$
quench the metastable, so the sink takes the sum; only the Penning branch
$f_{\rm P}=0.9$ leaves a lasting proton, and it appears as the source
$f_{\rm P}Q_{31}n_{2^3{\rm S}}n_{\rm HI}$ in the H$^+$ row. The remaining
10 percent is associative ionization to HeH$^+$, which the atomic systems
drop (in an atomic gas HeH$^+$ dissociatively recombines straight back) and
the molecular system carries explicitly.

\paragraph{The molecular rows.} \texttt{mol\_heh\_rows} adds four species
balances on the \citet{Koskinen2022} Table 1 network
(Sec.~\ref{sec:mol:network}). Schematically, writing $g$ for photorates and
$k_n$ for the Table-1 coefficients,
\begin{align}
\text{H}_2:\quad 0 &= k_6 n_e n_{\rm H_3^+} + k_9 n_{\rm H_2^+}n_{\rm HI}
   + k_{11}n_{\rm H_3^+}n_{\rm HI} + k_{15}n_{\rm HI}^{2}
   \nonumber\\ &\quad
   - \big[g_{\rm H_2} + g_{\rm LW} + (k_{10}+k_{13})n_{\rm HII}
          + k_{12}n_{\rm tot} + k_{14}n_e + k_8 n_{\rm H_2^+}
          + (k_{17}+k_{20}+k_{23})n_{\rm HeII} + k_{18}n_{\rm HeH^+}
          + k^{\rm ion}_{\rm H_2}n_{2^3{\rm S}}\big]n_{\rm H_2},
   \label{eq:ion:h2row}\\
\text{H}_2^+:\quad 0 &=
   \left(g_{\rm H_2}-g^{\rm di}-g^{\rm dd}-g^{\rm nd}\right)n_{\rm H_2}
   + k_{10}n_{\rm HII}n_{\rm H_2} + k_{11}n_{\rm H_3^+}n_{\rm HI}
   + k_{19}n_{\rm HeH^+}n_{\rm HI} + k_{23}n_{\rm HeII}n_{\rm H_2}
   \nonumber\\ &\quad
   + f_{\rm P}k^{\rm ion}_{\rm H_2}n_{2^3{\rm S}}n_{\rm H_2}
   - \left(k_5 n_e + k_8 n_{\rm H_2} + k_9 n_{\rm HI}\right)n_{\rm H_2^+},
   \label{eq:ion:h2prow}\\
\text{H}_3^+:\quad 0 &= k_8 n_{\rm H_2^+}n_{\rm H_2}
   + k_{13}n_{\rm HII}n_{\rm H_2} + k_{18}n_{\rm HeH^+}n_{\rm H_2}
   - \left[(k_6+k_7)n_e + k_{11}n_{\rm HI}\right]n_{\rm H_3^+},
   \label{eq:ion:h3prow}\\
\text{HeH}^+:\quad 0 &= k_{20}n_{\rm HeII}n_{\rm H_2}
   + (1-f_{\rm P})n_{2^3{\rm S}}
     \left(Q_{31}n_{\rm HI}+k^{\rm ion}_{\rm H_2}n_{\rm H_2}\right)
   - \left(k_{16}n_e+k_{18}n_{\rm H_2}+k_{19}n_{\rm HI}\right)n_{\rm HeH^+},
   \label{eq:ion:hehprow}
\end{align}
and the H$^+$ row gains
$(g^{\rm di}+2g^{\rm dd})n_{\rm H_2} + k_9 n_{\rm H_2^+}n_{\rm HI}
 + k_{17}n_{\rm HeII}n_{\rm H_2}
 - (k_{10}+k_{13})n_{\rm HII}n_{\rm H_2}$. The factor 2 on the double
photoionization branch is the stoichiometry, two protons out of one event,
while the H$_2$ destroyed is still one per event, which is why
\eqref{eq:ion:h2row} is not doubled with it.

\paragraph{The oxygen rows.} With the oxygen option
(Sec.~\ref{sec:mol:oxygen}; \texttt{thereis\_oxychem}, solved only by
\texttt{System\_HeH\_mol\_metals} since oxygen is a metal element) two more
unknowns join the cell system, $x_{\rm OH}=n_{\rm OH}/n_{\rm O}$ and
$x_{\rm H_2O}=n_{\rm H_2O}/n_{\rm O}$ with $n_{\rm O}$ the oxygen family of
the cell, and \texttt{oxygen\_carrier\_rows} adds their balances on the
reactions O1 (OH$+$H$_2\to$H$_2$O$+$H), O2 (O$+$H$_2\to$OH$+$H), their
reverses O1r and O2r, the three H$_2$O photolysis channels O3
($\to$OH$+$H), O4 ($\to$H$_2+$O($^1$D)) and O5 ($\to$O$+$H$+$H), and the OH
photolysis O7 ($\to$O$+$H):
\begin{align}
\text{OH}:\quad 0 &= k_{1r}n_{\rm H_2O}n_{\rm HI} + k_{2}n_{\rm O^0}n_{\rm H_2}
   + (J_3+J_4)\,n_{\rm H_2O}
   - k_{1}n_{\rm OH}n_{\rm H_2} - k_{2r}n_{\rm OH}n_{\rm HI} - J_7 n_{\rm OH},
   \label{eq:ion:ohrow}\\
\text{H}_2\text{O}:\quad 0 &= k_{1}n_{\rm OH}n_{\rm H_2}
   - k_{1r}n_{\rm H_2O}n_{\rm HI} - (J_3+J_4+J_5)\,n_{\rm H_2O},
   \label{eq:ion:h2orow}
\end{align}
and the H$_2$ row \eqref{eq:ion:h2row} gains the exchange with the water
cycle, $-(k_1 n_{\rm OH}+k_2 n_{\rm O^0})n_{\rm H_2}
+(k_{1r}n_{\rm H_2O}+k_{2r}n_{\rm OH})n_{\rm HI}$. Here $n_{\rm O^0}$ is the
\emph{free} atomic oxygen, the remainder of the family once the nuclei bound
in OH, H$_2$O and CO are removed, $n_{\rm O^0}=n_{\rm O}-n_{\rm OH}
-n_{\rm H_2O}-n_{\rm CO}$, which the metal block then splits over
O\,I/O\,II/O\,III from its own balance: the O\,I column of the output means
free atomic oxygen when the option is on. The hydrogen closure gains the
nuclei the carriers hold, $n_{\rm OH}+2n_{\rm H_2O}$. The forward
coefficients $k_1$ and $k_2$ are the \citet{Baulch2005} evaluations
($k_1=3.6\times10^{-16}\,T^{1.52}e^{-1740/T}$,
$k_2=6.34\times10^{-12}e^{-4000/T}+1.46\times10^{-9}e^{-9650/T}$
cm$^{3}$\,s$^{-1}$); the reverses are not transcribed but generated by
detailed balance, $k_{r}=k_{f}/K_c(T)$ with $K_c$ from the NIST-JANAF
Shomate Gibbs energies \citep{Chase1998} of H, H$_2$, O, OH, H$_2$O, CO and
C (\texttt{oxygen\_rates.f90}), so that the hot limit of the network is the
chemical equilibrium of the same thermochemistry. $J_3$, $J_4$, $J_5$ and
$J_7$ are the band-resolved photolysis rates of Sec.~\ref{sec:mol:oxygen}
summed over the four far-ultraviolet bands with their quantum yields. The
O4 branch appears in the OH row as if it made OH directly: its product
O($^1$D) has exactly one sink in the network, O6
(O($^1$D)$+$H$_2\to$OH$+$H), with a lifetime of $1.6\times10^{-3}$\,s at the
HD\,189733\,b base, so its local steady state is exact, the flux through O6
equals the O($^1$D) production $J_4 n_{\rm H_2O}$, and the pair O4$+$O6 is
one channel whose net effect on H$_2$ cancels (O4 makes one, O6 consumes
one), which is why only O1 and O2 enter the H$_2$ row; the O($^1$D) density
is evaluated for the output only. CO is \emph{not} a row of this system: it
is a transported carrier (Sec.~\ref{sec:mol:carrier}) with no formation
channel in the network, destroyed one-sidedly by
He$^+$$+$CO$\to$C$^+$$+$O$+$He at the Langevin value
$1.6\times10^{-9}$\,cm$^{3}$\,s$^{-1}$ \citep[UMIST RATE22 entry
4068;][]{Millar2024} and by photodissociation on the Lyman-Werner beam with
the \citet{Visser2009} shielding of the star-ward CO and H$_2$ columns, both
inside the carrier step; the cell system sees the CO of the cell as a frozen
background and takes the He$^+$ it consumes, $-k_{\rm D1}n_{\rm CO}n_{\rm
HeII}$, in the He$^+$ row, so the one reaction is complete over the pair of
operators and counted once in each. Each oxygen row is divided by its own
turnover rate $n_{\rm O}[(k_1+k_{1r}+k_2+k_{2r})n_{\rm H}+J_3+J_4+J_7]$ (OH)
and $n_{\rm O}[(k_1+k_{1r})n_{\rm H}+J_3+J_4+J_5]$ (H$_2$O) before the Newton
solve, as the molecular rows are, because the oxygen family is
$\sim10^{-4}$ of the hydrogen and the block would otherwise sit far below the
helium one.

\subsection{The rate coefficients}
\label{sec:ionization:rates}

All of the following live in \texttt{Cool\_coeff.f90} unless stated. The
default set is selected by \texttt{Legacy\_HHe\_rates: False} (default) and
\texttt{Atomic rate set: default}; the two alternatives are the
\citet{Hui1997} fits and the \citet{Koskinen2022} Table 1 fits, and the
three meet at the accessor functions at the end of that module.

\paragraph{Recombination.} Radiative recombination follows the \citet[2023 update]{Badnell2006}
total ground-level fit
\begin{equation}
\alpha_{\rm RR}(T) = \frac{A}
  {\sqrt{T/T_0}\left(1+\sqrt{T/T_0}\right)^{1-B'}
   \left(1+\sqrt{T/T_1}\right)^{1+B'}},
\qquad B' = B + C e^{-T_2/T},
\label{eq:ion:rrbadnell}
\end{equation}
implemented as \texttt{rr\_badnell} (its dummy \texttt{T0} keeps the
published notation and is not the code's temperature scale). Dielectronic
recombination of He\,{\sc ii} adds the adf09 three-term sum
$T^{-3/2}\sum_i c_i e^{-E_i/T}$. Case B is formed by subtracting the $n=1$
direct-capture level fit of \citet{Mao2016},
\begin{equation}
\alpha_1(T) = \frac{a_0\,\theta^{-b_0-c_0\ln\theta}
                    \left(1+a_2\theta^{-b_2}\right)}
                   {1+a_1\theta^{-b_1}},
\qquad \theta = k_{\rm B}T\ \mathrm{[eV]},
\label{eq:ion:rrmao}
\end{equation}
so that $\alpha_{\rm B} = \alpha_{\rm A}-\alpha_1$ for H\,{\sc ii},
He\,{\sc ii} and He\,{\sc iii}. The metal stages take Badnell RR $+$ adf09
DR through the same two forms, with the one gap (Ca\,{\sc ii}, K-like)
handled separately and noted at the code site.

\paragraph{Collisional ionization.} \citet{Voronov1997} Table 1,
\begin{equation}
\beta(T) = A\,\frac{\left(1+P\sqrt{U}\right)U^{K}e^{-U}}{X+U},
\qquad U = \frac{\Delta E}{k_{\rm B}T},
\label{eq:ion:voronov}
\end{equation}
with $(\Delta E,P,A,X,K)$ equal to $(13.6,0,2.91\times10^{-8},0.232,0.39)$
for H\,{\sc i}, $(24.6,0,1.75\times10^{-8},0.180,0.35)$ for He\,{\sc i} and
$(54.4,1,2.05\times10^{-9},0.265,0.25)$ for He\,{\sc ii}. The $\Delta E$
here is Voronov's own fit coefficient, not the measured potential; the
energy \emph{removed} per ionization is the measured one,
$E_{{\rm th}}$ in erg.

\paragraph{Charge exchange.} \texttt{charge\_exchange.f90} carries the 63
reactions of \citet{Huang2023} Table 4 plus one row that is not in it. The default active set is Group A, the 23 metal$+$H/H$^+$
reactions; \texttt{cx\_full} adds the metal$+$He (C) and metal$+$metal (D)
groups. The He$\leftrightarrow$H pair (Group B),
${\rm He}^0+{\rm H}^+ \rightleftharpoons {\rm He}^+ + {\rm H}^0$, is applied
separately by \texttt{he\_h\_cx\_fvec} so that it is available in
\emph{every} system containing helium and is never counted twice; it is
gated by \texttt{He\_H\_charge\_exchange} (default on). Group E is the one
reaction outside Table 4, ${\rm O}^{2+}+{\rm H}^0\to{\rm O}^{+}+{\rm H}^+$,
on by default at the published rate because its absence leaves the
O\,{\sc iii} profile wrong by decades in the region a transit probes;
\texttt{cx\_O2p\_H 0} in \texttt{metals.inp} restores the Table-4-only set
exactly. Each reaction enters as $R = k\,n_{\rm donor}n_{\rm acceptor}$, so
the zero-abundance guard is automatic.

\paragraph{Penning ionization of the metastable.} The total
He($2^3$S)$+$H ionization rate coefficient is the Maxwell-Boltzmann average
of the \citet{Cohen1971} and \citet{Morgner1979} cross sections as collected
by \citet{GarciaMunoz2025},
\begin{equation}
Q_{31}(T) = 10^{-9}\exp\!\left[-\frac{86.4804}{T}
   - 0.286766\ln T + 0.0868445\left(\ln T\right)^{2}
   - 5.73001\times10^{-3}\left(\ln T\right)^{3}\right]\ \mathrm{cm^{3}\,s^{-1}},
\label{eq:ion:penning}
\end{equation}
smooth, positive and single-peaked at $3.8\times10^{3}$\,K. The code notes
at that site why the two-branch power law of \citet{Taylor2025} Table 2 is
\emph{not} used: it jumps by a factor 1.5724 at 4000\,K,
which a rate coefficient cannot do, and it disagrees with that paper's own
Figure 19; the expression above agrees with Figure 19 to 10--20 percent.
The corresponding He($2^3$S)$+$H$_2$ coefficient is
$k = aT^{b}e^{c/T}$ with the \citet{Cohen1977} amplitudes of the
\citet{GarciaMunoz2025} network file, reproducing its Table A.5 to 0.13
percent. Both split 0.9\,:\,0.1 between the Penning and the associative
branch.

\subsection{Solving one cell}
\label{sec:ionization:solve}

\paragraph{The solver.} \texttt{solve\_ieq} (\texttt{newton\_solver.f90})
tries a damped Newton iteration with an \emph{analytic} Jacobian first
(\texttt{newton\_dense}: backtracking Armijo line search on
$\|f\|_2$, at most 60 iterations and 20 line-search steps), and on any
failure restores the entry guess and calls the in-tree MINPACK
\texttt{hybrd1} from the same guess, so the result can never be worse than
the forward-difference solve. Each residual routine ships its own exact
Jacobian; the photoionization rates are constants of the solve, so those
Jacobians are exact for the system actually solved
(Sec.~\ref{sec:radiation:rates}). Run-wide counters record how many cells each
route handled, and \texttt{EXHALE\_FORCE\_HYBRD1} forces the MINPACK path
for an A/B comparison. \texttt{use\_newton\_ieq} is on by default.

\paragraph{The sweep.} \texttt{ioniz\_eq} (\texttt{ionization\_equilibrium.f90})
takes the temperature and the total density as \emph{inputs} and overwrites
the species fractions in place; chemistry rearranges nucleons and electrons
and creates no mass, so the mass density is never corrected by this step and
the mass sum of the returned composition is a check against it
(\texttt{rho\_recon}), never a correction. The grid is traversed from the
top down in blocks; the columns of \eqref{eq:rad:column} are advanced with
the composition already solved, so the field each cell is solved at is
self-consistent with the states above it in one traversal. Inside a cell
the pair (field, composition) is iterated at most
\texttt{xuv\_self\_field\_passes} times, stopping when the accepted stage
fractions move by less than \texttt{xuv\_self\_field\_tol}. Cells within a
block are independent and the sweep is OpenMP-parallel with the per-cell
scratch declared \texttt{threadprivate}; blocks are of fixed even length so
that the vectorized transcendental calls pair the same lanes at any thread
count.

\paragraph{Acceptance classes.} Every cell's returned state carries a class,
and the contract is that classes 1, 2, 3 and 5 are \emph{roots} of the
requested equations (each passed the same two tests, a state inside the
element simplex and a normalized reaction residual at or below
\texttt{ieq\_res\_tol}) and differ only in how the root was found:

\begin{center}
\begin{tabular}{cl}
\toprule
class & meaning\\
\midrule
1 & root of the last attempt made\\
2 & root of the best earlier attempt\\
3 & root after projecting the closest attempt onto the element budget\\
4 & not a root; accepted under the relaxation amnesty and reported\\
5 & root of the constrained element-conserving continuation solve\\
6 & not a root, residual above \texttt{ieq\_nonroot\_res\_cap}: the cell
    \emph{keeps the composition it entered the sweep with}\\
\bottomrule
\end{tabular}
\end{center}

A residual that is not a finite number is a class-4 non-root by
construction, since NaN passes every tolerance test. A consumer counting
states its chemistry does not describe counts classes 4 and 6 and nothing
else.

\paragraph{The constrained continuation.}
\texttt{constrained\_chemical\_equilibrium.f90} is the second formulation of
the molecular network, and it removes the possibility of an unphysical
iterate instead of repairing one. The unknowns are $u_k = \ln n_k$, so
$n_k>0$ for every iterate, converged or not (the logarithmic composition
variables of the CEA method, \citealt{Gordon1994}); element
conservation is an explicit residual row per element rather than a closure,
so a species carrying nuclei of two elements (HeH$^+$, OH, H$_2$O) satisfies
both rows at the accepted root; and $n_e$ is the charge-neutrality sum, never
an unknown. Because the network is bistable in its starting point (a
molecular basin in the shielded base, an atomic basin above the H$_2\to$H
front), the solve is a natural-parameter continuation \citep{Allgower2003} in the
radiation field:
$\lambda$ scales the field from six decades below the cell's own, where the
composition is known in closed form from the H$_2$ dissociation equilibrium
of the local $p,T$, up to $\lambda=1$, each rung starting from the previous
rung's root and a failed rung being bisected geometrically. The rungs use the
general MINPACK \texttt{hybrd} and not \texttt{hybrd1}, whose fixed initial
trust radius $100\|x\|$ would open a first step of $e^{10^{4}}$ in every
density when $\|u\|\sim130$. A transported fraction, when one is imposed,
enters as a further constraint row.

\subsection{Ionization carried by the flow}
\label{sec:ionization:advected}

Where the recombination time is not short compared with the flow time, the
local balance of Sec.~\ref{sec:ionization:cellsystem} is not the answer. Two
places in the code carry the ionization state along the flow instead.

\paragraph{The implicit advection systems.} \texttt{System\_implicit\_adv\_H},
\texttt{\_HeH} and \texttt{\_HeH\_TR} solve, per hydrogen nucleus and with
$c_1 = \Delta r/v$ the residence time of the cell,
\begin{align}
0 &= x^{\rm up}_{\rm HI}-x_{\rm HI}
   + c_1\left[-\left(\Gamma_{\rm HI}+\beta_{\rm HI}x_e n_{\rm H}\right)x_{\rm HI}
              + \alpha_{\rm HII}x_{\rm HII}x_e n_{\rm H}\right],
   \label{eq:ion:advh}\\
0 &= x^{\rm up}_{\rm HeI}-x_{\rm HeI}
   + c_1\left[-\left(\Gamma_{\rm HeI}+\beta_{\rm HeI}x_e n_{\rm H}\right)x_{\rm HeI}
              + \alpha_{\rm HeII}x_{\rm HeII}x_e n_{\rm H}\right],
   \label{eq:ion:advhe}\\
0 &= x^{\rm up}_{\rm HeIII}-x_{\rm HeIII}
   + c_1\left[\left(\Gamma_{\rm HeII}+\beta_{\rm HeII}x_e n_{\rm H}\right)x_{\rm HeII}
              - \alpha_{\rm HeIII}x_{\rm HeIII}x_e n_{\rm H}\right],
   \label{eq:ion:advhe2}
\end{align}
with $x_e = x_{\rm HII} + ({\rm He/H})_{\rm local}(x_{\rm HeII}+2x_{\rm HeIII})
+ x_e^{\rm metal}$, and the same He$\leftrightarrow$H charge-exchange terms
on rows 1 and 2. The helium ratio is the global one on the legacy path and
the \emph{local}, radius-dependent one when the element diffusion is on,
since the global value would misstate $n_e$ by the separation factor. These
are the systems of the \texttt{\_adv} post-process
(Sec.~\ref{sec:postprocessing:adv}).

\paragraph{Transported protons.} With \texttt{Ionization transport} (default
off) H$^+$ is carried as a fifth transported species on the hydrodynamic
face mass flux inside the Runge-Kutta stages \eqref{eq:ssprk3}, in the same
way the molecular carriers are (Sec.~\ref{sec:mol:carrier}), and its local balance is then a constraint rather
than a closure.

\subsection{\texorpdfstring{Excited hydrogen and the Ly$\alpha$ field}{Excited hydrogen and the Lyman-alpha field}}
\label{sec:ionization:n2}

\texttt{excited\_hydrogen.f90} solves the $2s/2p$ statistical equilibrium of
\citet[their eqs.~(12)--(13)]{Christie2013} as a
$2\times2$ linear system,
\begin{equation}
\begin{pmatrix} L_{2p} & -M_{12}\\ -M_{21} & L_{2s}\end{pmatrix}
\begin{pmatrix} n_{2p}\\ n_{2s}\end{pmatrix}
=\begin{pmatrix} S_{2p}\\ S_{2s}\end{pmatrix},
\qquad
n_{2p} = \frac{S_{2p}L_{2s}+M_{12}S_{2s}}{L_{2p}L_{2s}-M_{12}M_{21}},
\label{eq:ion:n2pop}
\end{equation}
whose sources are the recombination cascade, collisional excitation from
$1s$ and radiative pumping by the Ly$\alpha$ mean intensity, and whose
sinks are radiative decay, the two-photon continuum, collisional
de-excitation, $\ell$-mixing and photoionization by the stellar Balmer
continuum. The atomic data are in one place
(\texttt{hydrogen\_n2\_rates.f90}): $A_{2p1s}=6.2649\times10^{8}$\,s$^{-1}$
(NIST ASD, the nonrelativistic dipole value carrying the reduced-mass
factor), $A_{2s1s}=8.26$\,s$^{-1}$, $g_{1s}=g_{2s}=2$, $g_{2p}=6$,
$f_{\rm Ly\alpha}=0.4162$, the \citeauthor{Christie2013} Table 2 collisional rates
such as
$q_{1s2p} = 1.71\times10^{-8}\,t_4^{-0.077}e^{-118400/T}$, their
detailed-balance inverses, and the \citet{Draine2011} case-B recombination and
its $2s/2p$ split,
\begin{equation}
\alpha_{\rm B} = 2.54\times10^{-13}\,
   t_4^{-0.8163-0.0208\ln t_4},
\qquad
\alpha_{2s} = \left(0.282+0.047t_4-0.006t_4^{2}\right)\alpha_{\rm B},
\label{eq:ion:alphaB}
\end{equation}
with $t_4 = T/10^{4}$\,K. The removal of a $2p$ atom without emission of a
line photon is
\begin{equation}
D_{2p} = n_e q_{2p1s} + \Gamma_{2p}
       + n_e C_{2p2s}\frac{A_{2s1s}}
         {A_{2s1s}+n_e(q_{2s1s}+C_{2s2p})+\Gamma_{2s}}.
\label{eq:ion:d2p}
\end{equation}

\paragraph{Balmer continuum.} The $n=2$ photoionization rate and the
photoelectric heating it delivers are integrals over
$[E_{\rm th}({\rm H},n{=}2), E_{\rm th}({\rm H})]$ of the run's own
spectrum with the hydrogenic cross section
$\sigma_2(E) = \sigma_2^{\rm thr}(E_2/E)^{3}$,
$\sigma_2^{\rm thr} = 1.4\times10^{-17}$\,cm$^{2}$:
\begin{equation}
\Gamma_2 = \frac{\sigma_2^{\rm thr}E_2^{3}}{h\,\nu_{\rm eV}}
           \int F_E\,E^{-4}\,{\rm d}E,
\qquad
h^{(1)}_2 = \sigma_2^{\rm thr}E_2^{3}
           \int F_E\left(E^{-3}-E_2E^{-4}\right){\rm d}E.
\label{eq:ion:balmer}
\end{equation}
The quadrature resolves the field it integrates, which a fixed rule cannot:
for a loaded table the nodes are the table's \emph{own} rows, and each
segment is integrated in closed form because the interpolation is a power
law there; for the power law the whole band is closed form; only for the
smooth Planck field is a composite five-point Gauss-Legendre rule on a
geometric subdivision used. The feedback is lagged-explicit: it is computed
once per timestep before the ionization and energy solve, and the frozen
arrays are read inside \texttt{ioniz\_eq} and never inside its Newton
iteration. The whole module is gated by \texttt{use\_excited\_H}, which is
armed by supplying \texttt{Stellar Teff [K]:} and
\texttt{Stellar radius [R\_sun]:}; it is \emph{off} by default.

\paragraph{The Ly$\alpha$ field.} \texttt{lya\_rt.f90} closes the line with
an escape probability rather than a Monte Carlo, because the columns
involved reach $\tau_0 \sim 10^{8}$ at line center. With
$\Delta\nu_{\rm D} = \nu_{\rm Ly\alpha}\sqrt{2kT/m_{\rm H}}/c$, the Voigt
parameter $a = (A_{2p1s}/4\pi)/\Delta\nu_{\rm D}$ and the line-center depth
accumulated cell by cell from $n_{\rm HI}C_{\rm Ly\alpha}/\Delta\nu_{\rm D}$,
the escape probability of a photon emitted at the cell center is
\begin{equation}
\beta = \frac{1}{1+\langle N\rangle},
\qquad
\langle N\rangle = \frac{4\sqrt6}{\pi^{2}}
   \frac{\tau_{\rm up}+\tau_{\rm dn}}{2}\,\Psi(\xi),
\qquad
\xi = \frac{\tau_{\rm dn}}{\tau_{\rm up}+\tau_{\rm dn}},
\qquad
\Psi(\xi) = \sum_{n\ \rm odd}\frac{\sin(n\pi\xi)}{n^{2}},
\label{eq:ion:neufeld}
\end{equation}
which is \citet{Neufeld1990} eq.~(3.27) at zero continuum
destruction, reducing for a mid-plane source to \citet{Harrington1973}
eq.~(40), $\langle N\rangle = 0.909316\,B$. The form $1/(1+\langle N\rangle)$ and not $1/\langle N\rangle$ is used
because $\langle N\rangle$ counts absorptions per photon generated, so a
photon is emitted $1+\langle N\rangle$ times before it leaves; this is exact
at both ends and needs no join between the thin and thick limits. Under
$R_{\rm II}$-A redistribution $\langle N\rangle$ is independent of $a$, so
the Voigt parameter survives only in the one-flight penetration of the
stellar beam. The trapped internal field and the stellar term are
\begin{equation}
J_{\rm int} = \frac{2h\nu^{3}}{c^{2}}\frac{g_{1s}}{g_{2p}}
   \frac{P\left(1-\beta\right)}{\left(A_{2p1s}\beta+D_{2p}\right)n_{\rm HI}},
\qquad
J_\star = \frac{\xi_{\rm day}F_{\rm Ly\alpha}}{4\pi\Delta\nu_\star}
          \left\langle T_s E\right\rangle,
\qquad
J_{\rm Ly\alpha} = J_{\rm int}+J_\star,
\label{eq:ion:jlya}
\end{equation}
with $P = \alpha_{\rm B}n_e n_{\rm HII} + C_{1s2p}n_e n_{\rm HI}$ the
internal production, $T_s = {\rm erfc}\!\left(x_1/\sqrt2 X_s\right)$ the
fraction of the broad stellar profile whose wings reach depth $\tau$, and
$\langle\cdot\rangle$ the mean of the beam over the cell's own optical
depth. The published solution requires $(a\tau)^{1/3}>10$
(\citeauthor{Neufeld1990} section Va); cells outside that are counted by
\texttt{lya\_wing\_domain\_record} and are carried there by the thin limit
$\beta\to1$, which \eqref{eq:ion:neufeld} reaches smoothly. The alternative
\texttt{jlya\_mode = 1} reads an externally computed $J_{\rm Ly\alpha}(r)$,
for instance from the LaRT Monte Carlo code, and interpolates it onto the
grid.
